% Options for packages loaded elsewhere
\PassOptionsToPackage{unicode}{hyperref}
\PassOptionsToPackage{hyphens}{url}
\documentclass[
  11pt,
]{article}
\usepackage{xcolor}
\usepackage[left=3cm,right=5cm,top=2cm,bottom=2cm]{geometry}
\usepackage{amsmath,amssymb}
\setcounter{secnumdepth}{5}
\usepackage{iftex}
\ifPDFTeX
  \usepackage[T1]{fontenc}
  \usepackage[utf8]{inputenc}
  \usepackage{textcomp} % provide euro and other symbols
\else % if luatex or xetex
  \usepackage{unicode-math} % this also loads fontspec
  \defaultfontfeatures{Scale=MatchLowercase}
  \defaultfontfeatures[\rmfamily]{Ligatures=TeX,Scale=1}
\fi
\usepackage{lmodern}
\ifPDFTeX\else
  % xetex/luatex font selection
\fi
% Use upquote if available, for straight quotes in verbatim environments
\IfFileExists{upquote.sty}{\usepackage{upquote}}{}
\IfFileExists{microtype.sty}{% use microtype if available
  \usepackage[]{microtype}
  \UseMicrotypeSet[protrusion]{basicmath} % disable protrusion for tt fonts
}{}
\makeatletter
\@ifundefined{KOMAClassName}{% if non-KOMA class
  \IfFileExists{parskip.sty}{%
    \usepackage{parskip}
  }{% else
    \setlength{\parindent}{0pt}
    \setlength{\parskip}{6pt plus 2pt minus 1pt}}
}{% if KOMA class
  \KOMAoptions{parskip=half}}
\makeatother
\usepackage{graphicx}
\makeatletter
\newsavebox\pandoc@box
\newcommand*\pandocbounded[1]{% scales image to fit in text height/width
  \sbox\pandoc@box{#1}%
  \Gscale@div\@tempa{\textheight}{\dimexpr\ht\pandoc@box+\dp\pandoc@box\relax}%
  \Gscale@div\@tempb{\linewidth}{\wd\pandoc@box}%
  \ifdim\@tempb\p@<\@tempa\p@\let\@tempa\@tempb\fi% select the smaller of both
  \ifdim\@tempa\p@<\p@\scalebox{\@tempa}{\usebox\pandoc@box}%
  \else\usebox{\pandoc@box}%
  \fi%
}
% Set default figure placement to htbp
\def\fps@figure{htbp}
\makeatother
% definitions for citeproc citations
\NewDocumentCommand\citeproctext{}{}
\NewDocumentCommand\citeproc{mm}{%
  \begingroup\def\citeproctext{#2}\cite{#1}\endgroup}
\makeatletter
 % allow citations to break across lines
 \let\@cite@ofmt\@firstofone
 % avoid brackets around text for \cite:
 \def\@biblabel#1{}
 \def\@cite#1#2{{#1\if@tempswa , #2\fi}}
\makeatother
\newlength{\cslhangindent}
\setlength{\cslhangindent}{1.5em}
\newlength{\csllabelwidth}
\setlength{\csllabelwidth}{3em}
\newenvironment{CSLReferences}[2] % #1 hanging-indent, #2 entry-spacing
 {\begin{list}{}{%
  \setlength{\itemindent}{0pt}
  \setlength{\leftmargin}{0pt}
  \setlength{\parsep}{0pt}
  % turn on hanging indent if param 1 is 1
  \ifodd #1
   \setlength{\leftmargin}{\cslhangindent}
   \setlength{\itemindent}{-1\cslhangindent}
  \fi
  % set entry spacing
  \setlength{\itemsep}{#2\baselineskip}}}
 {\end{list}}
\usepackage{calc}
\newcommand{\CSLBlock}[1]{\hfill\break\parbox[t]{\linewidth}{\strut\ignorespaces#1\strut}}
\newcommand{\CSLLeftMargin}[1]{\parbox[t]{\csllabelwidth}{\strut#1\strut}}
\newcommand{\CSLRightInline}[1]{\parbox[t]{\linewidth - \csllabelwidth}{\strut#1\strut}}
\newcommand{\CSLIndent}[1]{\hspace{\cslhangindent}#1}
\setlength{\emergencystretch}{3em} % prevent overfull lines
\providecommand{\tightlist}{%
  \setlength{\itemsep}{0pt}\setlength{\parskip}{0pt}}
\usepackage{amsmath,amssymb}
\usepackage{booktabs}
\usepackage{float}
\usepackage{booktabs}
\usepackage{longtable}
\usepackage{array}
\usepackage{multirow}
\usepackage{wrapfig}
\usepackage{float}
\usepackage{colortbl}
\usepackage{pdflscape}
\usepackage{tabu}
\usepackage{threeparttable}
\usepackage{threeparttablex}
\usepackage[normalem]{ulem}
\usepackage{makecell}
\usepackage{xcolor}
\usepackage{bookmark}
\IfFileExists{xurl.sty}{\usepackage{xurl}}{} % add URL line breaks if available
\urlstyle{same}
\hypersetup{
  pdftitle={Matched-Tuple Randomization for Factorial Clinical Trial Designs: Theory and Application to Trial-Ready Cohorts},
  pdfauthor={R.G. Thomas},
  hidelinks,
  pdfcreator={LaTeX via pandoc}}

\title{Matched-Tuple Randomization for Factorial Clinical Trial Designs:
Theory and Application to Trial-Ready Cohorts}
\author{R.G. Thomas}
\date{April 20, 2026}

\begin{document}
\maketitle
\begin{abstract}
Factorial clinical trial designs efficiently evaluate multiple treatment
factors simultaneously, but covariate balance across the resulting
treatment combinations is difficult to achieve through simple
randomization, particularly in moderate sample sizes. We propose
matched-tuple randomization, in which subjects are grouped into sets of
\(k\) individuals with similar baseline covariate profiles -- where
\(k\) equals the number of factorial treatment combinations -- and each
member of a tuple is assigned to a distinct treatment arm. We derive the
randomization distribution, variance estimators, and efficiency gains
relative to simple and stratified randomization. A simulation study
evaluates operating characteristics under scenarios motivated by
Alzheimer's disease prevention trials, where trial-ready cohorts provide
large pools of characterized subjects amenable to matching. Software
implementing the proposed methods is provided in the R package
\texttt{factorialrandom}.
\end{abstract}

{
\setcounter{tocdepth}{2}
\tableofcontents
}
\section{Introduction}\label{introduction}

\subsection{The Design Problem}\label{the-design-problem}

Clinical trials increasingly evaluate combinations of treatments.
Anti-amyloid plus anti-tau therapies in Alzheimer's disease, drug plus
behavioral intervention in oncology supportive care, and multi-component
public health programs all give rise to factorial treatment structures.
A \(2 \times 2\) factorial design assigns subjects to one of four
combinations: neither factor, factor A alone, factor B alone, or both.
The hidden replication
principle\textsuperscript{\citeproc{ref-fisherDesignExperiments1935}{1}}
ensures that each subject contributes to the estimation of both main
effects and their interaction, yielding the same precision as four
separate experiments of the same total size.

The challenge lies in covariate balance. With four or more treatment
arms, simple randomization can produce substantial imbalance on
prognostic covariates, inflating variance and threatening the
credibility of treatment comparisons. The problem is compounded in
moderate sample sizes (\(n = 50\) to \(500\) per arm) typical of Phase
II and prevention trials.

\subsection{Covariate Balance
Strategies}\label{covariate-balance-strategies}

Existing approaches to covariate balance -- stratified
randomization\textsuperscript{\citeproc{ref-kernanStratifiedRandomization1999}{2}},
minimization\textsuperscript{\citeproc{ref-pocockSequentialTreatment1975}{3}},
and
rerandomization\textsuperscript{\citeproc{ref-morganRerandomization2012}{4}}
-- were developed primarily for two-arm trials. Stratification is
limited to a few categorical covariates. Minimization extends to
multiple covariates but operates sequentially as subjects enroll, which
is unnecessary when the full pool is available at design time.
Rerandomization accepts only allocations meeting a balance criterion but
does not exploit the structure of factorial contrasts.

Matched-pair randomization represents the most aggressive form of
restricted
randomization.\textsuperscript{\citeproc{ref-baiOptimalityMatchedPairDesigns2022}{5}}
proved that among all stratified designs with equal allocation
probabilities, a specific matched-pair design achieves minimum
asymptotic
variance.\textsuperscript{\citeproc{ref-imaiVarianceIdentification2008}{6}}
and\textsuperscript{\citeproc{ref-fogartyRegressionAssisted2018}{7}}
developed the inferential framework. However, these results are
restricted to \(k = 2\) arms.

\subsection{Trial-Ready Cohorts}\label{trial-ready-cohorts}

The trial-ready cohort (TRC) paradigm, developed for Alzheimer's disease
prevention research, pre-screens and characterizes large pools of
potential trial participants. The TRC-PAD
project\textsuperscript{\citeproc{ref-aisenTRCPADOverview2020}{8},\citeproc{ref-aisenTRCPADExperience2020}{9}}
enrolled over 30,000 individuals in its web-based screening component
and characterized approximately 2,000 in-person participants with
cognitive testing, APOE genotyping, and amyloid biomarkers. The AHEAD
3-45 study\textsuperscript{\citeproc{ref-rafiiAHEAD2023}{10}} and
DIAN-TU
platform\textsuperscript{\citeproc{ref-batemanDIAN2012}{11},\citeproc{ref-bateman2017DIANTU}{12}}
draw on similar infrastructure.

These characterized pools create an opportunity: when the full cohort is
available before randomization, multivariate matching can be performed
optimally rather than sequentially. The rich baseline phenotyping
(amyloid PET, tau PET, cognitive composites, genetics) provides strong
prognostic covariates, and the factorial treatment structures emerging
in combination therapy trials create a natural demand for \(k > 2\)
matching.

\subsection{Contribution}\label{contribution}

We propose matched-tuple randomization for factorial designs. Subjects
are grouped into tuples of size \(k\) (where \(k\) is the number of
treatment combinations) based on multivariate baseline covariates, and
treatment assignment is constrained so that each tuple member receives a
distinct combination. We develop:

\begin{enumerate}
\def\labelenumi{\arabic{enumi}.}
\tightlist
\item
  Algorithms for forming matched \(k\)-tuples using Mahalanobis and
  prognostic score distances.
\item
  Exact and asymptotic variance formulas for factorial effects under the
  matched-tuple randomization distribution.
\item
  A simulation study evaluating efficiency gains under scenarios
  motivated by Alzheimer's disease prevention trials.
\item
  An R package implementing the methods.
\end{enumerate}

\section{Background}\label{background}

\subsection{Factorial Designs}\label{factorial-designs}

\subsubsection{Classical Framework}\label{classical-framework}

\textsuperscript{\citeproc{ref-fisherDesignExperiments1935}{1}}
introduced factorial experiments to evaluate multiple factors
simultaneously. Consider \(F\) factors, each at two levels (\(+\) and
\(-\)). The \(2^F\) factorial assigns subjects to all level
combinations. Main effects and interactions are defined as signed
averages of cell means. The key efficiency result is that every
observation contributes to every estimable effect.

\textsuperscript{\citeproc{ref-dasguptaCausalInference2015}{13}} recast
factorial designs within the potential outcomes framework. For a \(2^2\)
factorial, each subject has four potential outcomes
\(Y_i(00), Y_i(10), Y_i(01), Y_i(11)\) corresponding to the four
treatment combinations. The main effect of factor A is defined as:
\begin{equation}
  \tau_A = \frac{1}{N} \sum_{i=1}^{N}
    \left[ \frac{Y_i(10) + Y_i(11)}{2} -
           \frac{Y_i(00) + Y_i(01)}{2} \right]
\end{equation} and analogously for factor B and the interaction \(AB\).

\subsubsection{Factorial Designs in Clinical
Trials}\label{factorial-designs-in-clinical-trials}

Factorial designs are increasingly used in clinical trials to evaluate
combination therapies and multi-component
interventions.\textsuperscript{\citeproc{ref-collinsMultiphaseOptimization2007}{14}}
introduced the Multiphase Optimization Strategy (MOST), using factorial
experiments to identify active intervention components.

In Alzheimer's disease, the next generation of prevention trials is
expected to combine mechanisms of action: anti-amyloid antibodies with
anti-tau agents, secretase inhibitors with lifestyle interventions, or
multiple doses of the same agent with and without an adjunct
therapy\textsuperscript{\citeproc{ref-cummingsADTrialDesign2024}{15}}.
These combinations create natural factorial structures.

\subsection{Matched Randomization}\label{matched-randomization}

\subsubsection{\texorpdfstring{Matched-Pair Designs
(\(k = 2\))}{Matched-Pair Designs (k = 2)}}\label{matched-pair-designs-k-2}

In a matched-pair design, \(N\) subjects are grouped into \(M =
N/2\) pairs based on baseline covariates, and within each pair, one
subject is randomized to treatment and the other to control. The
treatment effect estimator is: \begin{equation}
  \hat{\tau} = \frac{1}{M} \sum_{m=1}^{M}
    (Y_{m,T} - Y_{m,C})
\end{equation} where \(Y_{m,T}\) and \(Y_{m,C}\) are the outcomes for
the treated and control subjects in pair \(m\).

Under the matched-pair randomization distribution, the variance of
\(\hat{\tau}\) is: \begin{equation}
  \text{Var}(\hat{\tau}) = \frac{1}{M^2} \sum_{m=1}^{M}
    \frac{(\tau_m - \bar{\tau})^2 +
          \text{Var}_m}{2}
\end{equation} where \(\tau_m = Y_{m}(1) - Y_{m}(0)\) is the within-pair
treatment effect and \(\text{Var}_m\) captures within-pair variability
in potential outcomes.

When pairs are well matched on prognostic covariates, the within-pair
treatment effects \(\tau_m\) are similar to each other, and the first
term is small. This is the source of precision gains.

\subsubsection{Optimal Matching}\label{optimal-matching}

\textsuperscript{\citeproc{ref-greevyOptimalMatching2004}{16}}
formulated optimal multivariate matching as a minimum-cost network flow
problem. For \(k = 2\), the problem reduces to minimum-weight perfect
matching on a bipartite graph, solvable in \(O(N^3)\) time. The distance
between subjects \(i\) and \(j\) is typically the Mahalanobis distance:
\begin{equation}
  d(i,j) = \sqrt{(X_i - X_j)' S^{-1} (X_i - X_j)}
\end{equation} where \(X_i\) is the covariate vector and \(S\) is the
sample covariance matrix.

\subsubsection{Rerandomization}\label{rerandomization}

\textsuperscript{\citeproc{ref-morganRerandomization2012}{4}} proposed
rerandomization as an intermediate approach between simple randomization
and deterministic matching. A randomization is accepted only if the
Mahalanobis distance between treatment group covariate means falls below
a threshold \(a\): \begin{equation}
  M = (\bar{X}_T - \bar{X}_C)' \hat{\text{cov}}(\bar{X}_T -
    \bar{X}_C)^{-1} (\bar{X}_T - \bar{X}_C) \le a
\end{equation}

\textsuperscript{\citeproc{ref-liRerandomizationAsymptotics2018}{17}}
derived the asymptotic distribution of treatment effect estimators under
rerandomization, showing it is a truncated normal with reduced variance.

\subsection{Trial-Ready Cohorts in Alzheimer's
Disease}\label{trial-ready-cohorts-in-alzheimers-disease}

\subsubsection{The Recruitment
Challenge}\label{the-recruitment-challenge}

Alzheimer's disease prevention trials face screening failure rates
exceeding 80\%, primarily because eligibility requires confirmation of
brain amyloid pathology via PET imaging or cerebrospinal fluid
analysis\textsuperscript{\citeproc{ref-sperlingTrialDesignAD2011}{18}}.
The cost of screening a single participant can exceed \$10,000. These
constraints motivate the trial-ready cohort approach, in which screening
and characterization are performed prospectively, creating a pool of
eligible participants ready for rapid enrollment.

\subsubsection{TRC-PAD}\label{trc-pad}

The Trial-Ready Cohort for Preclinical/Prodromal Alzheimer's Disease
(TRC-PAD) was established by the Alzheimer's Clinical Trials
Consortium\textsuperscript{\citeproc{ref-aisenTRCPADOverview2020}{8}}.
The architecture includes:

\begin{itemize}
\tightlist
\item
  An online longitudinal screening study (APT Webstudy) with quarterly
  cognitive assessments. Over 30,000 individuals consented by 2020.
\item
  Machine learning risk prediction using demographic, genetic, and
  cognitive data to identify individuals likely to have elevated brain
  amyloid.
\item
  In-person screening with cognitive battery, APOE genotyping, and
  plasma biomarkers (p-tau217, A\$\beta\$42/40
  ratio)\textsuperscript{\citeproc{ref-mcdonoughTRCScreening2024}{19}}.
\item
  Longitudinal follow-up of biomarker-confirmed eligible participants
  until a matching trial opens.
\end{itemize}

\subsubsection{AHEAD 3-45}\label{ahead-3-45}

The AHEAD 3-45 study\textsuperscript{\citeproc{ref-rafiiAHEAD2023}{10}}
tests lecanemab in preclinical AD using a sister-trial design: A3
(intermediate amyloid, Phase 2, imaging endpoints) and A45 (elevated
amyloid, Phase 3, cognitive endpoints). Both share a common screening
process drawing from TRC-PAD and direct recruitment. The study
demonstrates how TRC infrastructure feeds directly into clinical trials.

\subsubsection{DIAN-TU and EPAD}\label{dian-tu-and-epad}

The Dominantly Inherited Alzheimer
Network\textsuperscript{\citeproc{ref-batemanDIAN2012}{11}} maintains a
registry of autosomal dominant AD mutation carriers, providing a
genetically defined trial-ready cohort. The DIAN-TU
platform\textsuperscript{\citeproc{ref-bateman2017DIANTU}{12}}
implemented adaptive designs with shared placebo arms and cognitive
run-in periods.

The European Prevention of Alzheimer's Dementia (EPAD)
programme\textsuperscript{\citeproc{ref-solomonEPAD2022}{20}} recruited
over 2,000 participants into a longitudinal cohort across 39 European
institutions, explicitly designed as a readiness cohort for adaptive
proof-of-concept trials.

\subsubsection{Relevance to Matched Factorial
Designs}\label{relevance-to-matched-factorial-designs}

Trial-ready cohorts provide the essential prerequisites for
matched-tuple randomization: large characterized pools with rich
baseline phenotyping. The pool-to-sample ratio (available subjects
divided by trial enrollment target) determines matching quality.
TRC-PAD's pool of thousands of characterized individuals, combined with
the factorial treatment structures emerging in combination therapy
trials, creates an ideal application setting.

\section{Methods}\label{methods}

\subsection{Notation}\label{notation}

Consider \(N\) subjects to be randomized across a \(2^F\) factorial
design with \(k = 2^F\) treatment combinations. Subjects are grouped
into \(M = N/k\) matched tuples, each containing \(k\) subjects with
similar baseline covariate profiles. Within each tuple, the \(k\)
treatment combinations are assigned uniformly at random to the \(k\)
subjects.

Let \(X_i \in \mathbb{R}^p\) denote the baseline covariate vector for
subject \(i\). Let \(Y_i(z)\) denote the potential outcome for subject
\(i\) under treatment combination \(z \in
\{0, 1\}^F\). The observed outcome is \(Y_i = Y_i(Z_i)\) where \(Z_i\)
is the assigned treatment.

\subsection{Matching Algorithm}\label{matching-algorithm}

\subsubsection{Distance Metric}\label{distance-metric}

We consider two distance metrics for tuple formation:

\textbf{Mahalanobis distance.} For subjects \(i\) and \(j\):
\begin{equation}
  d_M(i,j) = \sqrt{(X_i - X_j)' S^{-1} (X_i - X_j)}
\end{equation} where \(S\) is the pooled sample covariance matrix of the
covariates.

\textbf{Prognostic score distance.} When historical data or pilot study
results are available, fit a model \(\hat{m}(X) =
E[Y(0) \mid X]\) predicting the outcome under the reference condition.
Then: \begin{equation}
  d_P(i,j) = |\hat{m}(X_i) - \hat{m}(X_j)|
\end{equation} Matching on the prognostic score concentrates precision
gains on the most relevant covariate dimension.

\subsubsection{Tuple Formation}\label{tuple-formation}

Given \(N\) subjects and tuple size \(k\), partition the subjects into
\(M = \lfloor N/k \rfloor\) tuples to minimize total within-tuple
dispersion: \begin{equation}
  \min_{\mathcal{P}} \sum_{m=1}^{M}
    \sum_{i < j \in \mathcal{T}_m} d(i,j)
\end{equation} where \(\mathcal{P} = \{\mathcal{T}_1, \ldots,
\mathcal{T}_M\}\) is a partition into tuples.

For \(k = 2\), this is optimal pair matching, solvable in polynomial
time. For \(k > 2\), we consider three approaches:

\begin{enumerate}
\def\labelenumi{\arabic{enumi}.}
\item
  \textbf{Sequential greedy.} Select a seed subject, find the \(k-1\)
  nearest neighbors, form a tuple, remove from pool, repeat. Order seeds
  by distance to the pool centroid (outermost first) to avoid stranding
  outliers.
\item
  \textbf{Hierarchical clustering.} Perform agglomerative clustering
  (Ward's method or complete linkage) and cut at the level producing
  \(M\) clusters of size \(k\). Post-process to enforce exact size \(k\)
  by splitting large clusters and merging small ones.
\item
  \textbf{Integer programming.} Formulate as an assignment problem:
  define binary variables \(x_{im} = 1\) if subject \(i\) is assigned to
  tuple \(m\), and solve the integer program with constraints
  \(\sum_m x_{im} = 1\) for all \(i\) and \(\sum_i x_{im} = k\) for all
  \(m\), minimizing total within-tuple distance. Feasible for
  \(N \lesssim 1000\) using standard solvers.
\end{enumerate}

\subsubsection{Pool Oversampling}\label{pool-oversampling}

When the trial-ready cohort provides \(N_{\text{pool}} > N\)
characterized subjects, matching quality improves. The algorithm selects
the \(N\) subjects that form the best \(M\) tuples, leaving
\(N_{\text{pool}} - N\) subjects unmatched. The pool-to-sample ratio
\(r = N_{\text{pool}} / N\) is a design parameter: larger \(r\) yields
tighter within-tuple matching at the cost of excluding more subjects.

\subsection{Randomization
Distribution}\label{randomization-distribution}

Within each tuple \(m\), the \(k\) treatment combinations are assigned
to the \(k\) subjects by drawing uniformly from the \(k!\) permutations.
Tuples are independent, so the randomization space has \((k!)^M\)
equally likely allocations.

This constrained randomization guarantees that for every tuple, each
treatment arm contains exactly one subject. Consequently, all factorial
contrasts (main effects, interactions) are perfectly balanced within
each tuple.

\subsection{Estimation}\label{estimation}

\subsubsection{Difference-in-Means}\label{difference-in-means}

For a factorial contrast \(c\) defined by a coefficient vector \(w_z\)
over treatment combinations (with \(\sum_z w_z = 0\)), the estimator is:
\begin{equation}
  \hat{\tau}_c = \frac{1}{M} \sum_{m=1}^{M}
    \sum_{z} w_z \, Y_{m}(z)
\end{equation} where \(Y_{m}(z)\) is the outcome of the subject in tuple
\(m\) assigned to treatment \(z\). This estimator is unbiased under the
matched-tuple randomization distribution.

\subsubsection{Variance Estimation}\label{variance-estimation}

The exact variance under matched-tuple randomization is:
\begin{equation}
  \text{Var}(\hat{\tau}_c) = \frac{1}{M^2}
    \sum_{m=1}^{M} V_m(c)
\end{equation} where \(V_m(c)\) depends on the within-tuple potential
outcomes. A conservative estimator replaces \(V_m(c)\) with the
within-tuple sample variance of the contrast-weighted outcomes.

\subsubsection{Regression Adjustment}\label{regression-adjustment}

Following\textsuperscript{\citeproc{ref-linAgnosticNotes2013}{21}}
and\textsuperscript{\citeproc{ref-fogartyRegressionAssisted2018}{7}}, we
augment the difference-in-means with a regression adjustment:
\begin{equation}
  \hat{\tau}_c^{\text{adj}} = \hat{\tau}_c -
    \hat{\beta}'(\bar{X}_c^{+} - \bar{X}_c^{-})
\end{equation} where \(\hat{\beta}\) is estimated from an outcome
regression and \(\bar{X}_c^{+}, \bar{X}_c^{-}\) are the covariate means
for the positive and negative arms of contrast \(c\). Under
matched-tuple randomization, \(\bar{X}_c^{+} - \bar{X}_c^{-}\) is
already close to zero, so the adjustment is small but further reduces
variance.

\subsection{Efficiency Analysis}\label{efficiency-analysis}

\subsubsection{Relative Efficiency}\label{relative-efficiency}

Define the relative efficiency of matched-tuple randomization versus
simple randomization as: \begin{equation}
  \text{RE} = \frac{\text{Var}_{\text{simple}}(\hat{\tau}_c)}
    {\text{Var}_{\text{matched}}(\hat{\tau}_c)}
\end{equation}

We derive RE as a function of the within-tuple intraclass correlation
\(\rho_w\) of the contrast-relevant potential outcomes. When tuples are
perfectly matched (\(\rho_w \to 1\)), \(\text{RE} \to k/(k-1) \cdot M\),
reflecting the elimination of between-tuple variation.

\subsubsection{Degrees of Freedom}\label{degrees-of-freedom}

A practical concern is the loss of degrees of freedom. Under simple
randomization with \(N\) subjects and \(k\) arms, the residual degrees
of freedom for variance estimation are \(N - k\). Under matched-tuple
randomization, the effective degrees of freedom are \(M - 1 = N/k - 1\).
When \(k\) is large and \(M\) is small, this loss can offset precision
gains from matching, particularly for coverage of confidence
intervals\textsuperscript{\citeproc{ref-martinMatchingPower1993}{22}}.

\section{Simulation Study}\label{simulation-study}

\subsection{Design}\label{design}

We evaluate five randomization strategies in a \(2 \times 2\) factorial
design (\(k = 4\)) under scenarios motivated by Alzheimer's disease
prevention trials.

\subsubsection{Data Generating
Mechanism}\label{data-generating-mechanism}

\begin{table}[!h]
\centering
\caption{\label{tab:dgm-params}Data generating mechanism parameters.}
\centering
\fontsize{10}{12}\selectfont
\begin{tabular}[t]{ll}
\toprule
Factor & Levels\\
\midrule
Total sample size ($N$) & 100, 200, 500\\
Covariate dimension ($p$) & 3, 8\\
Prognostic strength ($R^2$) & 0.1, 0.3, 0.5\\
Main effect A (SD units) & 0, 0.2, 0.5\\
Main effect B (SD units) & 0, 0.2, 0.5\\
\addlinespace
Interaction AB (SD units) & 0, 0.15\\
Pool/sample ratio ($r$) & 1.0, 1.5, 2.0, 3.0\\
Outcome distribution & Normal, skewed\\
\bottomrule
\end{tabular}
\end{table}

Subjects have \(p\) baseline covariates drawn from a multivariate normal
distribution with AR(1) correlation (\(\rho = 0.3\)). The outcome is
generated as: \begin{equation}
  Y_i = X_i' \beta + \tau_A \cdot A_i + \tau_B \cdot B_i +
    \tau_{AB} \cdot A_i B_i + \varepsilon_i
\end{equation} where \(\beta\) is calibrated to achieve the specified
\(R^2\), \(A_i\) and \(B_i\) are treatment indicators, and
\(\varepsilon_i \sim N(0, \sigma^2)\).

\subsubsection{Comparator Methods}\label{comparator-methods}

\begin{table}[!h]
\centering
\caption{\label{tab:comparators}Randomization strategies compared.}
\centering
\fontsize{10}{12}\selectfont
\begin{tabular}[t]{ll}
\toprule
Method & Description\\
\midrule
Simple randomization & Equal probability, no covariate balancing\\
Stratified randomization & Stratify on top 2 covariates (quartiles)\\
Minimization (Pocock-Simon) & Dynamic balancing on all covariates\\
Rerandomization (Mahalanobis) & Accept allocations with $M < \chi^2_{p,0.001}$\\
Matched-tuple (proposed) & Optimal $k$-tuples, random within-tuple\\
\bottomrule
\end{tabular}
\end{table}

\subsubsection{Evaluation Metrics}\label{evaluation-metrics}

For each scenario and method, we compute across 2,000 simulation
replicates:

\begin{itemize}
\tightlist
\item
  Covariate balance: maximum absolute standardized mean difference
  across all pairwise arm comparisons.
\item
  Bias of the main effect and interaction estimators.
\item
  Variance ratio (relative to simple randomization).
\item
  95\% confidence interval coverage.
\item
  Power (proportion of replicates rejecting \(H_0: \tau = 0\) at
  \(\alpha = 0.05\)).
\end{itemize}

\subsection{Results}\label{results}

\subsubsection{Covariate Balance}\label{covariate-balance}

\subsubsection{Efficiency Gains}\label{efficiency-gains}

\subsubsection{Coverage and Power}\label{coverage-and-power}

\subsubsection{Effect of Pool
Oversampling}\label{effect-of-pool-oversampling}

\section{Application: Alzheimer's Disease Prevention
Trial}\label{application-alzheimers-disease-prevention-trial}

We illustrate the proposed method using a design scenario motivated by
the AHEAD trial structure. Consider a \(2 \times
2\) factorial testing anti-amyloid therapy (yes/no) crossed with a
tau-targeting intervention (yes/no) in preclinical AD subjects drawn
from a trial-ready cohort.

\subsection{Design Parameters}\label{design-parameters}

\begin{itemize}
\tightlist
\item
  \textbf{Pool size:} 2,000 characterized TRC participants
\item
  \textbf{Trial enrollment:} 500 (125 per arm)
\item
  \textbf{Pool/sample ratio:} 4.0
\item
  \textbf{Matching covariates:} amyloid PET (Centiloids), tau PET
  (SUVR), APOE4 carrier status, baseline PACC score, age
\item
  \textbf{Primary endpoint:} Change in PACC at 4 years
\end{itemize}

\subsection{Matching Quality}\label{matching-quality}

\subsection{Treatment Effect
Estimation}\label{treatment-effect-estimation}

\section{Discussion}\label{discussion}

\subsection{Summary of Findings}\label{summary-of-findings}

\subsection{Practical Considerations}\label{practical-considerations}

\subsubsection{When to Use Matched-Tuple
Randomization}\label{when-to-use-matched-tuple-randomization}

Matched-tuple randomization is most advantageous when:

\begin{enumerate}
\def\labelenumi{\arabic{enumi}.}
\tightlist
\item
  The full pool of subjects is available before randomization (as in
  trial-ready cohorts), permitting global rather than sequential
  matching.
\item
  Strong prognostic covariates are measured at baseline.
\item
  The number of factorial arms \(k\) is not too large (\(k \le
  8\), i.e., up to \(2^3\) designs), to avoid excessive
  degrees-of-freedom loss.
\item
  The pool-to-sample ratio \(r\) exceeds 1, providing excess subjects
  for quality matching.
\end{enumerate}

\subsubsection{Limitations}\label{limitations}

The primary limitation is the loss of degrees of freedom relative to
unrestricted randomization. With \(M\) tuples, the effective degrees of
freedom for variance estimation are \(M - 1\), compared to \(N - k\)
under simple randomization. This concern is attenuated when \(N\) is
large enough that \(M = N/k\) remains sufficient (e.g., \(M > 30\)).

A second limitation is computational. Optimal \(k\)-tuple formation by
integer programming is feasible for \(N \lesssim
1000\) but may require heuristic approaches for larger pools. In the TRC
setting, pool sizes of 1,000 to 2,000 characterized subjects are
typical, placing the problem within the tractable range.

\subsubsection{Relation to
Rerandomization}\label{relation-to-rerandomization}

Rerandomization\textsuperscript{\citeproc{ref-morganRerandomization2012}{4}}
and matched-tuple randomization address the same goal -- covariate
balance -- through different mechanisms. Rerandomization restricts the
allocation space to a subset of all possible randomizations.
Matched-tuple randomization restricts both the allocation space and the
subject grouping. The two approaches can be combined: form tuples, then
rerandomize within tuples to further balance residual covariates, though
the marginal benefit of this combination may be small when tuples are
well matched.

\subsection{Future Work}\label{future-work}

Several extensions merit investigation:

\begin{itemize}
\tightlist
\item
  \textbf{Fractional factorials.} When \(2^F\) is too large, a
  \(2^{F-p}\) fractional factorial reduces \(k\). Matched-tuple
  randomization extends directly with \(k = 2^{F-p}\).
\item
  \textbf{Unequal allocation.} When factorial arms have unequal target
  sizes (e.g., shared placebo arm), the matching problem requires tuples
  of heterogeneous size.
\item
  \textbf{Adaptive enrichment.} Combining matched-tuple randomization
  with biomarker-driven adaptive enrichment designs.
\item
  \textbf{Time-to-event outcomes.} Extending the inference framework to
  survival endpoints.
\end{itemize}

\section{References}\label{references}

\subsection{Morris et al.~(2019) ADEMP
Compliance}\label{morris-et-al.-2019-ademp-compliance}

This simulation study was audited against the reporting standards
proposed by Morris, White, and
Crowther\textsuperscript{\citeproc{ref-morris2019using}{23}}. The full
audit is at \texttt{docs/morris-audit-2026-04-17.md}.

\textbf{Verdict:} Mostly compliant.

\textbf{Key gaps identified:}

\begin{itemize}
\tightlist
\item
  \texttt{n\_replicates\ =\ 2000} not justified via a Monte Carlo SE
  target.
\item
  MCSE columns are computed but not wired into the report tables ---
  TODO placeholders remain around \texttt{report.Rmd:605-649}.
\item
  \texttt{RNGkind(\textquotesingle{}L\textbackslash{}\textquotesingle{}Ecuyer-CMRG\textquotesingle{})}
  not pinned; no per-replicate \texttt{.Random.seed} snapshot stored.
\end{itemize}

A remediation plan is documented in the audit file and will be executed
in a subsequent revision.

\protect\phantomsection\label{refs}
\begin{CSLReferences}{0}{1}
\bibitem[\citeproctext]{ref-fisherDesignExperiments1935}
\CSLLeftMargin{1. }%
\CSLRightInline{Fisher RA. \emph{The Design of Experiments}. Oliver;
Boyd; 1935.}

\bibitem[\citeproctext]{ref-kernanStratifiedRandomization1999}
\CSLLeftMargin{2. }%
\CSLRightInline{Kernan WN, Viscoli CM, Makuch RW, Brass LM, Horwitz RI.
Stratified randomization for clinical trials. \emph{Journal of Clinical
Epidemiology}. 1999;52(1):19-26.
doi:\href{https://doi.org/10.1016/S0895-4356(98)00138-3}{10.1016/S0895-4356(98)00138-3}}

\bibitem[\citeproctext]{ref-pocockSequentialTreatment1975}
\CSLLeftMargin{3. }%
\CSLRightInline{Pocock SJ, Simon R. Sequential treatment assignment with
balancing for prognostic factors in the controlled clinical trial.
\emph{Biometrics}. 1975;31(1):103-115.
doi:\href{https://doi.org/10.2307/2529712}{10.2307/2529712}}

\bibitem[\citeproctext]{ref-morganRerandomization2012}
\CSLLeftMargin{4. }%
\CSLRightInline{Morgan KL, Rubin DB. Rerandomization to improve
covariate balance in experiments. \emph{The Annals of Statistics}.
2012;40(2):1263-1282.
doi:\href{https://doi.org/10.1214/12-AOS1008}{10.1214/12-AOS1008}}

\bibitem[\citeproctext]{ref-baiOptimalityMatchedPairDesigns2022}
\CSLLeftMargin{5. }%
\CSLRightInline{Bai Y. Optimality of {Matched-Pair Designs} in
{Randomized Controlled Trials}. \emph{American Economic Review}.
2022;112(12):3911-3940.
doi:\href{https://doi.org/10.1257/aer.20201856}{10.1257/aer.20201856}}

\bibitem[\citeproctext]{ref-imaiVarianceIdentification2008}
\CSLLeftMargin{6. }%
\CSLRightInline{Imai K. Variance identification and efficiency analysis
in randomized experiments under the matched-pair design.
\emph{Statistics in Medicine}. 2008;27(24):4857-4873.
doi:\href{https://doi.org/10.1002/sim.3337}{10.1002/sim.3337}}

\bibitem[\citeproctext]{ref-fogartyRegressionAssisted2018}
\CSLLeftMargin{7. }%
\CSLRightInline{Fogarty CB. Regression-assisted inference for the
average treatment effect in paired experiments. \emph{Biometrika}.
2018;105(4):994-1000.
doi:\href{https://doi.org/10.1093/biomet/asy034}{10.1093/biomet/asy034}}

\bibitem[\citeproctext]{ref-aisenTRCPADOverview2020}
\CSLLeftMargin{8. }%
\CSLRightInline{Aisen PS, Sperling RA, Cummings J, et al. The
{Trial-Ready Cohort} for {Preclinical}/{Prodromal Alzheimer}'s {Disease}
({TRC-PAD}) {Project}: {An Overview}. \emph{The Journal of Prevention of
Alzheimer's Disease}. 2020;7(4):208-212.
doi:\href{https://doi.org/10.14283/jpad.2020.45}{10.14283/jpad.2020.45}}

\bibitem[\citeproctext]{ref-aisenTRCPADExperience2020}
\CSLLeftMargin{9. }%
\CSLRightInline{Aisen PS, Jimenez-Maggiora GA, Rafii MS, Walter S, Raman
R. The {Trial-Ready Cohort} for {Preclinical} and {Prodromal
Alzheimer}'s {Disease} ({TRC-PAD}): {Experience} from the {First} 3
{Years}. \emph{The Journal of Prevention of Alzheimer's Disease}.
2020;7(4):234-241.
doi:\href{https://doi.org/10.14283/jpad.2020.47}{10.14283/jpad.2020.47}}

\bibitem[\citeproctext]{ref-rafiiAHEAD2023}
\CSLLeftMargin{10. }%
\CSLRightInline{Rafii MS, Sperling RA, Donohue MC, et al. The {AHEAD}
3--45 {Study}: {Design} of a prevention trial for {Alzheimer}'s disease.
\emph{Alzheimer's \& Dementia}. 2023;19(4):1227-1233.
doi:\href{https://doi.org/10.1002/alz.12748}{10.1002/alz.12748}}

\bibitem[\citeproctext]{ref-batemanDIAN2012}
\CSLLeftMargin{11. }%
\CSLRightInline{Bateman RJ, Xiong C, Benzinger TLS, et al. Clinical and
biomarker changes in dominantly inherited {Alzheimer}'s disease.
\emph{New England Journal of Medicine}. 2012;367(9):795-804.
doi:\href{https://doi.org/10.1056/NEJMoa1202753}{10.1056/NEJMoa1202753}}

\bibitem[\citeproctext]{ref-bateman2017DIANTU}
\CSLLeftMargin{12. }%
\CSLRightInline{Bateman RJ, Benzinger TLS, Berry S, et al. The {DIAN-TU
Next Generation Alzheimer}'s prevention trial: Adaptive design and
disease progression model. \emph{Alzheimer's \& Dementia}.
2017;13(1):8-19.
doi:\href{https://doi.org/10.1016/j.jalz.2016.07.005}{10.1016/j.jalz.2016.07.005}}

\bibitem[\citeproctext]{ref-dasguptaCausalInference2015}
\CSLLeftMargin{13. }%
\CSLRightInline{Dasgupta T, Pillai NS, Rubin DB. Causal inference from
2\textsuperscript{K} factorial designs by using potential outcomes.
\emph{Journal of the Royal Statistical Society: Series B (Statistical
Methodology)}. 2015;77(4):727-753.
doi:\href{https://doi.org/10.1111/rssb.12085}{10.1111/rssb.12085}}

\bibitem[\citeproctext]{ref-collinsMultiphaseOptimization2007}
\CSLLeftMargin{14. }%
\CSLRightInline{Collins LM, Murphy SA, Strecher V. The multiphase
optimization strategy ({MOST}) and the sequential multiple assignment
randomized trial ({SMART}): New methods for more potent {eHealth}
interventions. \emph{American Journal of Preventive Medicine}.
2007;32(5):S112-S118.
doi:\href{https://doi.org/10.1016/j.amepre.2007.01.022}{10.1016/j.amepre.2007.01.022}}

\bibitem[\citeproctext]{ref-cummingsADTrialDesign2024}
\CSLLeftMargin{15. }%
\CSLRightInline{Cummings J, Zhou Y, Lee G, Zhong K, Fonseca J, Cheng F.
Alzheimer's disease drug development pipeline: 2024. \emph{Alzheimer's
\& Dementia: Translational Research \& Clinical Interventions}.
2024;10(2):e12465.
doi:\href{https://doi.org/10.1002/trc2.12465}{10.1002/trc2.12465}}

\bibitem[\citeproctext]{ref-greevyOptimalMatching2004}
\CSLLeftMargin{16. }%
\CSLRightInline{Greevy R, Lu B, Silber JH, Rosenbaum P. Optimal
multivariate matching before randomization. \emph{Biostatistics}.
2004;5(2):263-275.
doi:\href{https://doi.org/10.1093/biostatistics/5.2.263}{10.1093/biostatistics/5.2.263}}

\bibitem[\citeproctext]{ref-liRerandomizationAsymptotics2018}
\CSLLeftMargin{17. }%
\CSLRightInline{Li X, Ding P, Rubin DB. Asymptotic theory of
rerandomization in treatment-control experiments. \emph{Proceedings of
the National Academy of Sciences}. 2018;115(37):9157-9162.
doi:\href{https://doi.org/10.1073/pnas.1808191115}{10.1073/pnas.1808191115}}

\bibitem[\citeproctext]{ref-sperlingTrialDesignAD2011}
\CSLLeftMargin{18. }%
\CSLRightInline{Sperling RA, Aisen PS, Beckett LA, et al. Toward
defining the preclinical stages of {Alzheimer}'s disease:
Recommendations from the {National Institute} on {Aging} --
{Alzheimer}'s {Association} workgroups on diagnostic guidelines for
{Alzheimer}'s disease. \emph{Alzheimer's \& Dementia}.
2011;7(3):280-292.
doi:\href{https://doi.org/10.1016/j.jalz.2011.03.003}{10.1016/j.jalz.2011.03.003}}

\bibitem[\citeproctext]{ref-mcdonoughTRCScreening2024}
\CSLLeftMargin{19. }%
\CSLRightInline{McDonough IM, Nosheny RL, Sperling RA, Aisen PS, Raman
R. Plasma biomarker prescreening for a {Trial-Ready Cohort} for
{Alzheimer}'s prevention trials. \emph{The Journal of Prevention of
Alzheimer's Disease}. 2024;11(3):608-617.
doi:\href{https://doi.org/10.14283/jpad.2024.26}{10.14283/jpad.2024.26}}

\bibitem[\citeproctext]{ref-solomonEPAD2022}
\CSLLeftMargin{20. }%
\CSLRightInline{Solomon A, Kivipelto M, Molinuevo JL, Tom B, Ritchie CW.
The {European Prevention} of {Alzheimer}'s {Dementia Programme}: {An
Innovative Medicines Initiative-funded} partnership to facilitate
secondary prevention of {Alzheimer}'s disease dementia. \emph{Frontiers
in Neurology}. 2022;13:1051543.
doi:\href{https://doi.org/10.3389/fneur.2022.1051543}{10.3389/fneur.2022.1051543}}

\bibitem[\citeproctext]{ref-linAgnosticNotes2013}
\CSLLeftMargin{21. }%
\CSLRightInline{Lin W. Agnostic notes on regression adjustments to
experimental data: Reexamining {Freedman's} critique. \emph{The Annals
of Applied Statistics}. 2013;7(1):295-318.
doi:\href{https://doi.org/10.1214/12-AOAS583}{10.1214/12-AOAS583}}

\bibitem[\citeproctext]{ref-martinMatchingPower1993}
\CSLLeftMargin{22. }%
\CSLRightInline{Martin DC, Diehr P, Perrin EB, Koepsell TD. The effect
of matching on the power of randomized community intervention studies.
\emph{Statistics in Medicine}. 1993;12(3--4):329-338.
doi:\href{https://doi.org/10.1002/sim.4780120315}{10.1002/sim.4780120315}}

\bibitem[\citeproctext]{ref-morris2019using}
\CSLLeftMargin{23. }%
\CSLRightInline{Morris TP, White IR, Crowther MJ. Using simulation
studies to evaluate statistical methods. \emph{Statistics in Medicine}.
2019;38(11):2074-2102.
doi:\href{https://doi.org/10.1002/sim.8086}{10.1002/sim.8086}}

\end{CSLReferences}

\end{document}
